\documentclass[11pt,a4paper]{ctexart}
\usepackage[margin=2.2cm]{geometry}
\usepackage{hyperref}
\usepackage{longtable}
\usepackage{enumitem}
\setlist{nosep}
\hypersetup{colorlinks=true,linkcolor=blue,urlcolor=blue}

\title{LeetCode 前 100 题伪代码整理}
\author{}
\date{\today}

\begin{document}
\maketitle
\tableofcontents
\newpage

\section*{说明}
\addcontentsline{toc}{section}{说明}
本文按 LeetCode 题号 1--100 整理常见解法伪代码，用于复习算法思路。伪代码侧重关键流程和复杂度，不包含完整可提交代码。

\newcommand{\problem}[3]{
\subsection{#1. #2}
\textbf{思路：} #3
}

\problem{1}{Two Sum}{哈希表保存已遍历数字的位置，遇到目标补数时直接返回。}
\begin{verbatim}
map = empty hash map
for i from 0 to n-1:
    need = target - nums[i]
    if need in map: return [map[need], i]
    map[nums[i]] = i
\end{verbatim}

\problem{2}{Add Two Numbers}{按链表位逐位相加，维护进位并生成新链表。}
\begin{verbatim}
carry = 0; dummy = new node
while l1 or l2 or carry:
    sum = carry + value(l1) + value(l2)
    append node(sum mod 10)
    carry = sum div 10
    move l1 and l2 if possible
return dummy.next
\end{verbatim}

\problem{3}{Longest Substring Without Repeating Characters}{滑动窗口维护无重复区间，哈希表记录字符最后位置。}
\begin{verbatim}
left = 0; ans = 0; last = empty map
for right from 0 to len(s)-1:
    c = s[right]
    if c in last and last[c] >= left: left = last[c] + 1
    last[c] = right
    ans = max(ans, right - left + 1)
return ans
\end{verbatim}

\problem{4}{Median of Two Sorted Arrays}{二分较短数组的切分位置，使左右两侧数量相等且有序。}
\begin{verbatim}
ensure A is shorter than B
binary search i in [0, len(A)]:
    j = (m+n+1)/2 - i
    if Aleft <= Bright and Bleft <= Aright:
        if odd: return max(Aleft, Bleft)
        else: return (max(lefts)+min(rights))/2
    else if Aleft > Bright: search left
    else: search right
\end{verbatim}

\problem{5}{Longest Palindromic Substring}{枚举每个中心，分别扩展奇数和偶数长度回文。}
\begin{verbatim}
best = ""
for center from 0 to n-1:
    update best by expand(center, center)
    update best by expand(center, center+1)
return best

expand(l, r):
    while l>=0 and r<n and s[l]==s[r]: l--, r++
    return s[l+1 ... r-1]
\end{verbatim}

\problem{6}{Zigzag Conversion}{按行模拟字符移动方向，到顶或到底时反向。}
\begin{verbatim}
if numRows == 1: return s
rows = array of empty strings
row = 0; dir = 1
for c in s:
    rows[row] += c
    if row == 0: dir = 1
    if row == numRows-1: dir = -1
    row += dir
return concatenate rows
\end{verbatim}

\problem{7}{Reverse Integer}{逐位弹出数字并检查 32 位整数边界。}
\begin{verbatim}
ans = 0
while x != 0:
    digit = x mod 10
    x = truncate(x / 10)
    if ans will overflow after ans*10+digit: return 0
    ans = ans*10 + digit
return ans
\end{verbatim}

\problem{8}{String to Integer (atoi)}{跳过空格，读取符号和连续数字，过程中截断到边界。}
\begin{verbatim}
i = first non-space index
sign = read optional + or -
ans = 0
while i<n and s[i] is digit:
    ans = ans*10 + digit
    if sign*ans exceeds bound: return bound
    i++
return sign * ans
\end{verbatim}

\problem{9}{Palindrome Number}{负数不是回文；反转后一半数字与前一半比较。}
\begin{verbatim}
if x < 0 or (x % 10 == 0 and x != 0): return false
rev = 0
while x > rev:
    rev = rev*10 + x%10
    x = x/10
return x == rev or x == rev/10
\end{verbatim}

\problem{10}{Regular Expression Matching}{动态规划，dp[i][j] 表示 s 前 i 个与 p 前 j 个是否匹配。}
\begin{verbatim}
dp[0][0] = true
for j where p[j-1] == '*': dp[0][j] = dp[0][j-2]
for i from 1..m:
  for j from 1..n:
    if p[j-1] matches s[i-1]: dp[i][j] = dp[i-1][j-1]
    else if p[j-1] == '*':
        dp[i][j] = dp[i][j-2]
        if p[j-2] matches s[i-1]: dp[i][j] |= dp[i-1][j]
return dp[m][n]
\end{verbatim}

\problem{11}{Container With Most Water}{双指针从两端向内移动，每次移动较短边。}
\begin{verbatim}
l = 0; r = n-1; ans = 0
while l < r:
    ans = max(ans, min(h[l],h[r])*(r-l))
    if h[l] < h[r]: l++ else: r--
return ans
\end{verbatim}

\problem{12}{Integer to Roman}{从大到小贪心匹配罗马数值和符号。}
\begin{verbatim}
pairs = [(1000,"M"),(900,"CM"),...,(1,"I")]
ans = ""
for value, symbol in pairs:
    while num >= value:
        ans += symbol
        num -= value
return ans
\end{verbatim}

\problem{13}{Roman to Integer}{若当前字符小于右侧字符则减，否则加。}
\begin{verbatim}
ans = 0
for i from 0 to n-1:
    v = value(s[i])
    if i+1<n and v < value(s[i+1]): ans -= v
    else: ans += v
return ans
\end{verbatim}

\problem{14}{Longest Common Prefix}{用第一个字符串作候选前缀，逐个缩短直到匹配所有字符串。}
\begin{verbatim}
prefix = strs[0]
for each str in strs:
    while str does not start with prefix:
        remove last char from prefix
        if prefix empty: return ""
return prefix
\end{verbatim}

\problem{15}{3Sum}{排序后枚举第一个数，内层双指针找相反数并跳过重复。}
\begin{verbatim}
sort nums; ans = []
for i from 0 to n-3:
    if i>0 and nums[i]==nums[i-1]: continue
    l=i+1; r=n-1
    while l<r:
        sum = nums[i]+nums[l]+nums[r]
        if sum==0: record triple; skip duplicate l,r
        else if sum<0: l++
        else: r--
return ans
\end{verbatim}

\problem{16}{3Sum Closest}{排序后固定一个数，双指针更新最接近 target 的和。}
\begin{verbatim}
sort nums; best = nums[0]+nums[1]+nums[2]
for i from 0 to n-3:
    l=i+1; r=n-1
    while l<r:
        sum = nums[i]+nums[l]+nums[r]
        if abs(sum-target) < abs(best-target): best = sum
        if sum < target: l++ else if sum > target: r-- else return sum
return best
\end{verbatim}

\problem{17}{Letter Combinations of a Phone Number}{回溯构造每一位数字对应的字母组合。}
\begin{verbatim}
if digits empty: return []
backtrack(pos, path):
    if pos == len(digits): add path; return
    for c in letters[digits[pos]]:
        backtrack(pos+1, path+c)
return all paths
\end{verbatim}

\problem{18}{4Sum}{排序后枚举前两个数，后两个数用双指针，注意去重和剪枝。}
\begin{verbatim}
sort nums; ans=[]
for i from 0..n-4:
  skip duplicate i
  for j from i+1..n-3:
    skip duplicate j
    l=j+1; r=n-1
    while l<r:
      sum = nums[i]+nums[j]+nums[l]+nums[r]
      if sum==target: record; skip duplicates
      else if sum<target: l++ else: r--
return ans
\end{verbatim}

\problem{19}{Remove Nth Node From End of List}{哑节点配合快慢指针，快指针先走 n+1 步。}
\begin{verbatim}
dummy.next = head
fast = slow = dummy
repeat n+1 times: fast = fast.next
while fast:
    fast = fast.next; slow = slow.next
slow.next = slow.next.next
return dummy.next
\end{verbatim}

\problem{20}{Valid Parentheses}{栈保存左括号，遇到右括号时检查栈顶是否匹配。}
\begin{verbatim}
stack = []
for c in s:
    if c is left bracket: push c
    else:
        if stack empty or top not matching c: return false
        pop
return stack empty
\end{verbatim}

\problem{21}{Merge Two Sorted Lists}{双指针归并两个有序链表。}
\begin{verbatim}
dummy = new node; cur = dummy
while l1 and l2:
    append smaller node to cur
append remaining non-empty list
return dummy.next
\end{verbatim}

\problem{22}{Generate Parentheses}{回溯保证左括号数量不超过 n，右括号不超过左括号。}
\begin{verbatim}
dfs(left, right, path):
    if len(path)==2*n: add path; return
    if left < n: dfs(left+1, right, path+"(")
    if right < left: dfs(left, right+1, path+")")
\end{verbatim}

\problem{23}{Merge k Sorted Lists}{小根堆维护每个链表当前最小节点。}
\begin{verbatim}
push head of each non-empty list into min heap
dummy = new node
while heap not empty:
    node = pop min
    append node
    if node.next: push node.next
return dummy.next
\end{verbatim}

\problem{24}{Swap Nodes in Pairs}{迭代处理每两个节点，调整前驱指针。}
\begin{verbatim}
dummy.next = head; pre = dummy
while pre.next and pre.next.next:
    a = pre.next; b = a.next
    pre.next = b; a.next = b.next; b.next = a
    pre = a
return dummy.next
\end{verbatim}

\problem{25}{Reverse Nodes in k-Group}{每次检查是否有 k 个节点，有则局部反转并连接。}
\begin{verbatim}
dummy.next=head; groupPrev=dummy
loop:
    kth = kth node after groupPrev
    if kth is null: break
    groupNext = kth.next
    reverse [groupPrev.next, kth]
    oldHead = groupPrev.next
    groupPrev.next = kth
    oldHead.next = groupNext
    groupPrev = oldHead
return dummy.next
\end{verbatim}

\problem{26}{Remove Duplicates from Sorted Array}{慢指针指向去重后数组末尾，快指针扫描新值。}
\begin{verbatim}
if nums empty: return 0
slow = 1
for fast from 1 to n-1:
    if nums[fast] != nums[fast-1]:
        nums[slow] = nums[fast]; slow++
return slow
\end{verbatim}

\problem{27}{Remove Element}{快慢指针保留所有不等于 val 的元素。}
\begin{verbatim}
slow = 0
for x in nums:
    if x != val:
        nums[slow] = x; slow++
return slow
\end{verbatim}

\problem{28}{Find the Index of the First Occurrence in a String}{KMP 构造前缀函数后在线匹配。}
\begin{verbatim}
build lps array for needle
j = 0
for i from 0 to len(haystack)-1:
    while j>0 and haystack[i]!=needle[j]: j = lps[j-1]
    if equal: j++
    if j == len(needle): return i-j+1
return -1
\end{verbatim}

\problem{29}{Divide Two Integers}{转为负数避免溢出，用倍增减法求商。}
\begin{verbatim}
handle overflow case
make dividend and divisor negative; sign = result sign
ans = 0
while dividend <= divisor:
    value = divisor; count = 1
    while value+value >= dividend and value+value < 0:
        value += value; count += count
    dividend -= value; ans += count
return sign * ans
\end{verbatim}

\problem{30}{Substring with Concatenation of All Words}{按单词长度分组滑动窗口，维护窗口内单词计数。}
\begin{verbatim}
need = count of words
for offset in 0..wordLen-1:
    left = offset; window = empty count
    for right step wordLen:
        word = s[right:right+wordLen]
        add word; shrink while window[word] > need[word]
        if window has all words: record left
return ans
\end{verbatim}

\problem{31}{Next Permutation}{从右找第一个下降点，交换其右侧更大的最小元素，再反转后缀。}
\begin{verbatim}
i = n-2
while i>=0 and nums[i]>=nums[i+1]: i--
if i>=0:
    j = n-1
    while nums[j] <= nums[i]: j--
    swap nums[i], nums[j]
reverse nums[i+1 ... n-1]
\end{verbatim}

\problem{32}{Longest Valid Parentheses}{栈保存下标，栈底是当前无效边界。}
\begin{verbatim}
stack = [-1]; ans = 0
for i from 0 to n-1:
    if s[i]=='(': push i
    else:
        pop
        if stack empty: push i
        else: ans = max(ans, i - top)
return ans
\end{verbatim}

\problem{33}{Search in Rotated Sorted Array}{二分判断哪一半有序，再决定目标在哪边。}
\begin{verbatim}
l=0; r=n-1
while l<=r:
    mid=(l+r)/2
    if nums[mid]==target: return mid
    if nums[l] <= nums[mid]:
        if nums[l]<=target<nums[mid]: r=mid-1 else: l=mid+1
    else:
        if nums[mid]<target<=nums[r]: l=mid+1 else: r=mid-1
return -1
\end{verbatim}

\problem{34}{Find First and Last Position of Element in Sorted Array}{两次二分分别找第一个大于等于 target 和大于 target 的位置。}
\begin{verbatim}
left = lower_bound(nums, target)
right = lower_bound(nums, target+epsilon) - 1
if left <= right: return [left, right]
else return [-1, -1]
\end{verbatim}

\problem{35}{Search Insert Position}{标准 lower bound。}
\begin{verbatim}
l=0; r=n
while l<r:
    mid=(l+r)/2
    if nums[mid] < target: l=mid+1
    else: r=mid
return l
\end{verbatim}

\problem{36}{Valid Sudoku}{用集合检查每行、每列、每个 3x3 宫是否重复。}
\begin{verbatim}
rows, cols, boxes = arrays of sets
for each cell (i,j):
    if board[i][j] == '.': continue
    b = (i/3)*3 + j/3
    if val in rows[i] or cols[j] or boxes[b]: return false
    add val to rows[i], cols[j], boxes[b]
return true
\end{verbatim}

\problem{37}{Sudoku Solver}{回溯填空格，使用行列宫集合快速判断合法性。}
\begin{verbatim}
collect empty cells
dfs(k):
    if k == number of empties: return true
    choose cell k
    for digit 1..9:
        if digit valid in row, col, box:
            place digit
            if dfs(k+1): return true
            remove digit
    return false
\end{verbatim}

\problem{38}{Count and Say}{从 "1" 开始，逐轮压缩连续相同字符。}
\begin{verbatim}
s = "1"
repeat n-1 times:
    next = ""
    scan groups of equal chars in s:
        next += count + char
    s = next
return s
\end{verbatim}

\problem{39}{Combination Sum}{排序后回溯，同一数字可重复使用。}
\begin{verbatim}
sort candidates
dfs(start, remain, path):
    if remain == 0: add path; return
    for i from start to end:
        if candidates[i] > remain: break
        dfs(i, remain-candidates[i], path+[candidates[i]])
\end{verbatim}

\problem{40}{Combination Sum II}{排序后回溯，每个数字最多用一次，同层跳过重复。}
\begin{verbatim}
sort candidates
dfs(start, remain, path):
    if remain == 0: add path; return
    for i from start to end:
        if i>start and candidates[i]==candidates[i-1]: continue
        if candidates[i] > remain: break
        dfs(i+1, remain-candidates[i], path+[candidates[i]])
\end{verbatim}

\problem{41}{First Missing Positive}{原地把数字 x 放到下标 x-1，最后找第一个不匹配位置。}
\begin{verbatim}
for i from 0 to n-1:
    while 1<=nums[i]<=n and nums[nums[i]-1] != nums[i]:
        swap nums[i], nums[nums[i]-1]
for i from 0 to n-1:
    if nums[i] != i+1: return i+1
return n+1
\end{verbatim}

\problem{42}{Trapping Rain Water}{双指针维护左右最高墙，较低一侧决定当前蓄水。}
\begin{verbatim}
l=0; r=n-1; leftMax=0; rightMax=0; ans=0
while l<r:
    if height[l] < height[r]:
        leftMax=max(leftMax,height[l]); ans+=leftMax-height[l]; l++
    else:
        rightMax=max(rightMax,height[r]); ans+=rightMax-height[r]; r--
return ans
\end{verbatim}

\problem{43}{Multiply Strings}{模拟竖式乘法，结果数组长度最多 m+n。}
\begin{verbatim}
res = array of zeros length m+n
for i from m-1 downto 0:
  for j from n-1 downto 0:
    mul = digit(num1[i]) * digit(num2[j])
    sum = mul + res[i+j+1]
    res[i+j+1] = sum % 10
    res[i+j] += sum / 10
return res without leading zeros
\end{verbatim}

\problem{44}{Wildcard Matching}{动态规划处理 ? 和 *，星号可匹配空串或继续消耗字符。}
\begin{verbatim}
dp[0][0]=true
for j where p[j-1]=='*': dp[0][j]=dp[0][j-1]
for i=1..m:
  for j=1..n:
    if p[j-1]=='?' or equal: dp[i][j]=dp[i-1][j-1]
    else if p[j-1]=='*': dp[i][j]=dp[i][j-1] or dp[i-1][j]
return dp[m][n]
\end{verbatim}

\problem{45}{Jump Game II}{贪心按层扫描，当前层结束时必须跳一次并扩展最远边界。}
\begin{verbatim}
steps=0; end=0; far=0
for i from 0 to n-2:
    far = max(far, i+nums[i])
    if i == end:
        steps++; end = far
return steps
\end{verbatim}

\problem{46}{Permutations}{回溯，每次选择一个未使用元素。}
\begin{verbatim}
dfs(path):
    if len(path)==n: add path; return
    for i in 0..n-1:
        if used[i]: continue
        used[i]=true; dfs(path+[nums[i]]); used[i]=false
\end{verbatim}

\problem{47}{Permutations II}{排序后回溯，若前一个相同元素未使用则跳过当前元素。}
\begin{verbatim}
sort nums
dfs(path):
    if len(path)==n: add path; return
    for i in 0..n-1:
        if used[i]: continue
        if i>0 and nums[i]==nums[i-1] and not used[i-1]: continue
        choose nums[i]; dfs; unchoose
\end{verbatim}

\problem{48}{Rotate Image}{先沿主对角线转置，再水平翻转每一行。}
\begin{verbatim}
for i in 0..n-1:
  for j in i+1..n-1: swap matrix[i][j], matrix[j][i]
for each row:
  reverse row
\end{verbatim}

\problem{49}{Group Anagrams}{把排序后的字符串作为 key 分组。}
\begin{verbatim}
map = empty dictionary
for str in strs:
    key = sort characters of str
    map[key].append(str)
return values of map
\end{verbatim}

\problem{50}{Pow(x, n)}{快速幂，指数二进制拆分；负指数取倒数。}
\begin{verbatim}
if n < 0: x = 1/x; n = -n
ans = 1
while n > 0:
    if n odd: ans *= x
    x *= x
    n = n / 2
return ans
\end{verbatim}

\problem{51}{N-Queens}{按行回溯，集合记录列、主对角线、副对角线占用。}
\begin{verbatim}
dfs(row):
    if row == n: add board; return
    for col in 0..n-1:
        if col or row-col or row+col occupied: continue
        place queen; mark; dfs(row+1); unmark
\end{verbatim}

\problem{52}{N-Queens II}{与 N 皇后相同，只统计方案数，不保存棋盘。}
\begin{verbatim}
count = 0
dfs(row):
    if row == n: count++; return
    for each legal col:
        mark; dfs(row+1); unmark
return count
\end{verbatim}

\problem{53}{Maximum Subarray}{Kadane 算法：当前位置最大子数组要么延续，要么从当前数重开。}
\begin{verbatim}
cur = best = nums[0]
for i from 1 to n-1:
    cur = max(nums[i], cur + nums[i])
    best = max(best, cur)
return best
\end{verbatim}

\problem{54}{Spiral Matrix}{维护上下左右边界，按右、下、左、上顺序收缩。}
\begin{verbatim}
top=0; bottom=m-1; left=0; right=n-1
while top<=bottom and left<=right:
    traverse top row; top++
    traverse right col; right--
    if top<=bottom: traverse bottom row; bottom--
    if left<=right: traverse left col; left++
\end{verbatim}

\problem{55}{Jump Game}{贪心维护能到达的最远位置，扫描中若当前位置超过最远则失败。}
\begin{verbatim}
far = 0
for i from 0 to n-1:
    if i > far: return false
    far = max(far, i + nums[i])
return true
\end{verbatim}

\problem{56}{Merge Intervals}{按起点排序，依次合并与当前区间重叠的区间。}
\begin{verbatim}
sort intervals by start
ans = []
for interval in intervals:
    if ans empty or interval.start > ans.last.end:
        append interval
    else:
        ans.last.end = max(ans.last.end, interval.end)
return ans
\end{verbatim}

\problem{57}{Insert Interval}{先加入新区间左侧，再合并所有重叠区间，最后加入右侧。}
\begin{verbatim}
append intervals ending before new.start
while current starts before or at new.end:
    new.start = min(new.start, current.start)
    new.end = max(new.end, current.end)
append merged new
append remaining intervals
\end{verbatim}

\problem{58}{Length of Last Word}{从字符串末尾跳过空格，再统计连续非空格长度。}
\begin{verbatim}
i = n-1
while i>=0 and s[i]==' ': i--
len = 0
while i>=0 and s[i]!=' ': len++, i--
return len
\end{verbatim}

\problem{59}{Spiral Matrix II}{按螺旋顺序向矩阵填入 1 到 n^2。}
\begin{verbatim}
num=1; top=0; bottom=n-1; left=0; right=n-1
while num <= n*n:
    fill top row; top++
    fill right col; right--
    fill bottom row; bottom--
    fill left col; left++
return matrix
\end{verbatim}

\problem{60}{Permutation Sequence}{用阶乘数系统逐位确定第 k 个排列。}
\begin{verbatim}
k = k-1; nums = [1..n]
for pos from n downto 1:
    block = factorial(pos-1)
    index = k / block
    append nums[index]; remove it
    k = k % block
return answer
\end{verbatim}

\problem{61}{Rotate List}{链表成环后在新尾位置断开。}
\begin{verbatim}
if head null: return head
compute length len and tail
k = k % len
if k == 0: return head
tail.next = head
steps to newTail = len-k-1
newHead = newTail.next
newTail.next = null
return newHead
\end{verbatim}

\problem{62}{Unique Paths}{二维或一维 DP，dp[j] 表示到当前格子的路径数。}
\begin{verbatim}
dp = array of 1 length n
for i from 1 to m-1:
    for j from 1 to n-1:
        dp[j] += dp[j-1]
return dp[n-1]
\end{verbatim}

\problem{63}{Unique Paths II}{障碍格路径数为 0，其余格等于上方加左方。}
\begin{verbatim}
dp[0] = 1 if start not obstacle else 0
for each row i:
    for each col j:
        if obstacle[i][j]: dp[j] = 0
        else if j>0: dp[j] += dp[j-1]
return dp[n-1]
\end{verbatim}

\problem{64}{Minimum Path Sum}{原地或一维 DP，当前位置最小和来自上方或左方较小者。}
\begin{verbatim}
for each cell (i,j):
    if start: continue
    up = dp value from above or infinity
    left = dp value from left or infinity
    grid[i][j] += min(up, left)
return grid[m-1][n-1]
\end{verbatim}

\problem{65}{Valid Number}{按有限状态或规则扫描符号、小数点、指数和数字。}
\begin{verbatim}
seenDigit=false; seenDot=false; seenExp=false
for each char c with index i:
    if digit: seenDigit=true
    else if c in '+-': valid only at start or after e/E
    else if c=='.': valid only if no dot and no exp
    else if c in 'eE': valid only if no exp and seenDigit; reset seenDigit=false
    else return false
return seenDigit
\end{verbatim}

\problem{66}{Plus One}{从末尾加一，处理连续进位。}
\begin{verbatim}
for i from n-1 downto 0:
    if digits[i] < 9:
        digits[i]++; return digits
    digits[i] = 0
return [1] + digits
\end{verbatim}

\problem{67}{Add Binary}{从两个字符串末尾开始按位相加并处理进位。}
\begin{verbatim}
i=m-1; j=n-1; carry=0; ans=""
while i>=0 or j>=0 or carry:
    sum = carry + bit(a,i) + bit(b,j)
    prepend sum%2 to ans
    carry = sum/2
    i--; j--
return ans
\end{verbatim}

\problem{68}{Text Justification}{贪心取每行最多单词，再按空格规则分配。}
\begin{verbatim}
while words remain:
    take maximal words fitting maxWidth
    if last line or one word:
        left justify with single spaces then pad right
    else:
        spaces = maxWidth - total letters
        distribute spaces evenly; earlier gaps get extras
return lines
\end{verbatim}

\problem{69}{Sqrt(x)}{二分查找最大整数 y，使 y*y <= x。}
\begin{verbatim}
l=0; r=x; ans=0
while l<=r:
    mid=(l+r)/2
    if mid <= x/mid: ans=mid; l=mid+1
    else: r=mid-1
return ans
\end{verbatim}

\problem{70}{Climbing Stairs}{斐波那契 DP，f[i]=f[i-1]+f[i-2]。}
\begin{verbatim}
if n<=2: return n
a=1; b=2
for i from 3 to n:
    c=a+b; a=b; b=c
return b
\end{verbatim}

\problem{71}{Simplify Path}{按斜杠分割，栈处理目录名、点和双点。}
\begin{verbatim}
stack = []
for token in split(path, '/'):
    if token empty or token == ".": continue
    if token == "..": pop if stack not empty
    else: push token
return "/" + join(stack, "/")
\end{verbatim}

\problem{72}{Edit Distance}{经典 DP，插入、删除、替换三种转移。}
\begin{verbatim}
dp[i][0]=i; dp[0][j]=j
for i=1..m:
  for j=1..n:
    if word1[i-1]==word2[j-1]: dp[i][j]=dp[i-1][j-1]
    else dp[i][j]=1+min(dp[i-1][j], dp[i][j-1], dp[i-1][j-1])
return dp[m][n]
\end{verbatim}

\problem{73}{Set Matrix Zeroes}{用第一行第一列作为标记，额外记录第一列是否清零。}
\begin{verbatim}
firstColZero = any matrix[i][0] == 0
for cells excluding first row/col:
    if cell==0: mark row and col in first column/row
zero cells according to marks
if matrix[0][0]==0: zero first row
if firstColZero: zero first col
\end{verbatim}

\problem{74}{Search a 2D Matrix}{把矩阵看作长度 m*n 的有序数组进行二分。}
\begin{verbatim}
l=0; r=m*n-1
while l<=r:
    mid=(l+r)/2
    val = matrix[mid/n][mid%n]
    if val==target: return true
    else if val<target: l=mid+1
    else: r=mid-1
return false
\end{verbatim}

\problem{75}{Sort Colors}{荷兰国旗算法，三个区域分别放 0、1、2。}
\begin{verbatim}
zero=0; i=0; two=n-1
while i<=two:
    if nums[i]==0: swap nums[zero],nums[i]; zero++; i++
    else if nums[i]==2: swap nums[i],nums[two]; two--
    else: i++
\end{verbatim}

\problem{76}{Minimum Window Substring}{滑动窗口满足需求后尽量收缩左边界。}
\begin{verbatim}
need = char counts of t; missing = len(t)
left=0; best=none
for right in 0..len(s)-1:
    if need[s[right]] > 0: missing--
    need[s[right]]--
    while missing == 0:
        update best with [left,right]
        need[s[left]]++
        if need[s[left]] > 0: missing++
        left++
return best substring or ""
\end{verbatim}

\problem{77}{Combinations}{回溯从递增数字中选择 k 个，可用剩余数量剪枝。}
\begin{verbatim}
dfs(start, path):
    if len(path)==k: add path; return
    for x from start to n - (k-len(path)) + 1:
        dfs(x+1, path+[x])
\end{verbatim}

\problem{78}{Subsets}{逐个元素选择或不选择，或递归枚举起点。}
\begin{verbatim}
ans = []
dfs(start, path):
    add path to ans
    for i from start to n-1:
        dfs(i+1, path+[nums[i]])
return ans
\end{verbatim}

\problem{79}{Word Search}{从每个格子开始 DFS，访问过的格子临时标记。}
\begin{verbatim}
for each cell:
    if dfs(cell, 0): return true

dfs(i,j,k):
    if k == len(word): return true
    if out of bound or visited or board[i][j]!=word[k]: return false
    mark visited
    search four neighbors with k+1
    unmark
\end{verbatim}

\problem{80}{Remove Duplicates from Sorted Array II}{保留每个数字最多两次，慢指针比较前两位。}
\begin{verbatim}
slow = 0
for x in nums:
    if slow < 2 or x != nums[slow-2]:
        nums[slow] = x; slow++
return slow
\end{verbatim}

\problem{81}{Search in Rotated Sorted Array II}{带重复值的旋转数组二分，无法判断时收缩左边界。}
\begin{verbatim}
l=0; r=n-1
while l<=r:
    mid=(l+r)/2
    if nums[mid]==target: return true
    if nums[l]==nums[mid]: l++; continue
    if nums[l] < nums[mid]:
        decide target in left sorted half
    else:
        decide target in right sorted half
return false
\end{verbatim}

\problem{82}{Remove Duplicates from Sorted List II}{哑节点加前驱指针，整段跳过重复值。}
\begin{verbatim}
dummy.next=head; pre=dummy; cur=head
while cur:
    if cur.next and cur.val==cur.next.val:
        val=cur.val
        while cur and cur.val==val: cur=cur.next
        pre.next=cur
    else:
        pre=cur; cur=cur.next
return dummy.next
\end{verbatim}

\problem{83}{Remove Duplicates from Sorted List}{遍历有序链表，跳过与当前节点值相同的后继。}
\begin{verbatim}
cur=head
while cur and cur.next:
    if cur.val == cur.next.val: cur.next = cur.next.next
    else: cur = cur.next
return head
\end{verbatim}

\problem{84}{Largest Rectangle in Histogram}{单调递增栈，弹出时以该高度计算最大矩形。}
\begin{verbatim}
append sentinel 0
stack = [-1]; ans=0
for i from 0 to n:
    while stack top valid and heights[top] > heights[i]:
        h = heights[pop]
        w = i - stack.top - 1
        ans = max(ans, h*w)
    push i
return ans
\end{verbatim}

\problem{85}{Maximal Rectangle}{逐行把 1 的连续高度转成柱状图，调用最大矩形算法。}
\begin{verbatim}
heights = zeros of column count
for each row:
    for each col:
        heights[col] = heights[col]+1 if cell=='1' else 0
    ans = max(ans, largestRectangle(heights))
return ans
\end{verbatim}

\problem{86}{Partition List}{维护小于 x 和大于等于 x 两条链，最后拼接。}
\begin{verbatim}
smallDummy, bigDummy
for each node:
    if node.val < x: append to small
    else: append to big
terminate big list
small.next = bigDummy.next
return smallDummy.next
\end{verbatim}

\problem{87}{Scramble String}{区间 DP 或记忆化搜索，枚举切分并判断交换/不交换两种情况。}
\begin{verbatim}
memo dfs(aStart, bStart, len):
    if substrings equal: return true
    if char counts differ: return false
    for k from 1 to len-1:
        if dfs(a,b,k) and dfs(a+k,b+k,len-k): return true
        if dfs(a,b+len-k,k) and dfs(a+k,b,len-k): return true
    return false
\end{verbatim}

\problem{88}{Merge Sorted Array}{从两个数组有效部分末尾向前归并，避免覆盖。}
\begin{verbatim}
i=m-1; j=n-1; k=m+n-1
while j>=0:
    if i>=0 and nums1[i] > nums2[j]:
        nums1[k]=nums1[i]; i--
    else:
        nums1[k]=nums2[j]; j--
    k--
\end{verbatim}

\problem{89}{Gray Code}{第 i 个格雷码为 i xor (i >> 1)。}
\begin{verbatim}
ans=[]
for i from 0 to 2^n - 1:
    ans.append(i xor (i >> 1))
return ans
\end{verbatim}

\problem{90}{Subsets II}{排序后回溯，同层跳过重复元素。}
\begin{verbatim}
sort nums
dfs(start, path):
    add path
    for i from start to n-1:
        if i>start and nums[i]==nums[i-1]: continue
        dfs(i+1, path+[nums[i]])
\end{verbatim}

\problem{91}{Decode Ways}{DP，单字符 1-9 可解码，双字符 10-26 可解码。}
\begin{verbatim}
dp[0]=1
for i from 1 to n:
    if s[i-1] != '0': dp[i] += dp[i-1]
    if i>=2 and 10 <= int(s[i-2:i]) <= 26: dp[i] += dp[i-2]
return dp[n]
\end{verbatim}

\problem{92}{Reverse Linked List II}{找到 left 前驱，头插法反转指定区间。}
\begin{verbatim}
dummy.next=head; pre=dummy
move pre to node before left
cur = pre.next
repeat right-left times:
    next = cur.next
    cur.next = next.next
    next.next = pre.next
    pre.next = next
return dummy.next
\end{verbatim}

\problem{93}{Restore IP Addresses}{回溯选择 4 段，每段长度 1 到 3 且数值合法。}
\begin{verbatim}
dfs(pos, parts):
    if parts count == 4:
        if pos == len(s): add join(parts,'.')
        return
    for len in 1..3:
        seg = s[pos:pos+len]
        if seg valid (no leading zero unless single, <=255):
            dfs(pos+len, parts+[seg])
\end{verbatim}

\problem{94}{Binary Tree Inorder Traversal}{迭代栈：一路向左入栈，弹出访问后转右子树。}
\begin{verbatim}
stack=[]; cur=root; ans=[]
while cur or stack:
    while cur: push cur; cur=cur.left
    cur=pop; ans.append(cur.val)
    cur=cur.right
return ans
\end{verbatim}

\problem{95}{Unique Binary Search Trees II}{递归枚举根节点，组合所有左右子树。}
\begin{verbatim}
build(l, r):
    if l > r: return [null]
    ans=[]
    for rootVal from l to r:
        for left in build(l, rootVal-1):
            for right in build(rootVal+1, r):
                ans.append(new Tree(rootVal,left,right))
    return ans
return build(1,n)
\end{verbatim}

\problem{96}{Unique Binary Search Trees}{Catalan DP，枚举根节点划分左右子树数量。}
\begin{verbatim}
dp[0]=dp[1]=1
for nodes from 2 to n:
    for leftSize from 0 to nodes-1:
        dp[nodes] += dp[leftSize] * dp[nodes-1-leftSize]
return dp[n]
\end{verbatim}

\problem{97}{Interleaving String}{二维 DP，判断 s3 前 i+j 个能否由 s1 前 i 个和 s2 前 j 个交错组成。}
\begin{verbatim}
if len(s1)+len(s2)!=len(s3): return false
dp[0][0]=true
for i=0..m:
  for j=0..n:
    if i>0 and s1[i-1]==s3[i+j-1]: dp[i][j] |= dp[i-1][j]
    if j>0 and s2[j-1]==s3[i+j-1]: dp[i][j] |= dp[i][j-1]
return dp[m][n]
\end{verbatim}

\problem{98}{Validate Binary Search Tree}{递归传递当前节点允许的开区间上下界。}
\begin{verbatim}
valid(node, low, high):
    if node null: return true
    if node.val <= low or node.val >= high: return false
    return valid(node.left, low, node.val) and valid(node.right, node.val, high)
return valid(root, -inf, +inf)
\end{verbatim}

\problem{99}{Recover Binary Search Tree}{中序遍历应递增，找到两个逆序节点并交换值。}
\begin{verbatim}
prev=null; first=null; second=null
inorder(node):
    inorder(node.left)
    if prev and prev.val > node.val:
        if first is null: first = prev
        second = node
    prev = node
    inorder(node.right)
inorder(root)
swap first.val and second.val
\end{verbatim}

\problem{100}{Same Tree}{递归比较两个节点是否同时为空、值相等且左右子树相同。}
\begin{verbatim}
same(p, q):
    if p null and q null: return true
    if p null or q null: return false
    if p.val != q.val: return false
    return same(p.left,q.left) and same(p.right,q.right)
\end{verbatim}

\end{document}
